\section{Langevin Dynamics：基于 Score 的采样}

\subsection{核心思想}

前面我们建立了 score function 的概念，并看到它编码了关于概率分布的丰富信息。现在的关键问题是：如果我们知道某个分布的 score function，能否从该分布中采样？

答案是肯定的，Langevin dynamics（朗之万动力学）正是实现这一目标的方法。它源自统计物理学，核心思想是：通过 score function 引导的迭代过程，可以从\textbf{任意初始分布}出发，最终收敛到目标分布。

\begin{importantbox}{Langevin Dynamics 采样过程}
给定目标分布 $q(x)$ 的 score function $\nabla_x \log q(x)$，Langevin dynamics 的迭代更新公式为：
$$x_{k+1} = x_k + \eta\, \nabla_x \log q(x_k) + \sqrt{2\eta}\, z_k, \quad z_k \sim \mathcal{N}(0, I)$$
其中 $\eta > 0$ 是步长参数。当 $\eta \to 0$ 且迭代次数 $K \to \infty$ 时，$x_K$ 的分布收敛到 $q(x)$。
\end{importantbox}

\begin{knowledgebox}{Langevin Dynamics 的三个组成部分}
Langevin dynamics 的更新公式由三部分组成：
\begin{enumerate}
    \item \textbf{当前位置} $x_k$：粒子的当前状态。
    \item \textbf{Score 引导} $\eta\, \nabla_x \log q(x_k)$：将粒子推向高概率密度区域的确定性力。
    \item \textbf{随机噪声} $\sqrt{2\eta}\, z_k$：防止粒子陷入局部极值，确保对目标分布的充分探索。
\end{enumerate}
如果只有 score 引导而没有噪声注入，粒子只会收敛到概率密度的极大值点（模式），而不会形成正确的分布。噪声项的存在确保了最终样本的多样性。
\end{knowledgebox}

\subsection{Langevin Dynamics 的关键性质}

Langevin dynamics 有一个非凡的性质：它\textbf{不需要知道}概率密度函数 $q(x)$ 本身，只需要 score function $\nabla_x \log q(x)$。

\begin{importantbox}{Score Function 足以完全描述分布}
Score function $\nabla_x \log q(x)$ 在满足一定正则条件时，唯一确定了概率分布 $q(x)$（up to normalization constant）。这是因为：
\begin{itemize}
    \item $\nabla_x \log q(x) = \frac{\nabla_x q(x)}{q(x)}$，包含了 $q(x)$ 的局部梯度信息。
    \item 通过 Langevin dynamics，我们可以从 score function 恢复出采样过程。
    \item 这也意味着：即使不知道概率密度函数的归一化常数（这在高维空间中通常难以计算），只要知道 score function 就足以进行采样。
\end{itemize}
\end{importantbox}

在传统的概率模型中，从复杂分布采样通常需要计算累积分布函数的逆（inverse CDF），这在高维空间中几乎不可行。Langevin dynamics 提供了一种替代方案：只需计算对数概率密度对数据点的梯度，然后通过迭代过程即可采样。

\begin{warningbox}{Langevin Dynamics 的实践限制}
理论上的收敛保证需要步长 $\eta \to 0$ 和无限次迭代，实践中两者都无法满足。有限步长和有限迭代次数会引入近似误差。此外，在多模态分布中，Langevin dynamics 可能需要很长时间才能在不同模式之间跳转（mixing problem）。
\end{warningbox}

\subsection{本章小结}

Langevin dynamics 是一种基于 score function 的采样方法，它通过``score 引导 + 随机噪声注入''的迭代过程，可以从任意初始分布出发收敛到目标分布。这个方法的核心优势在于：只需要 score function 而非概率密度函数本身就可以进行采样，从而避免了高维空间中归一化常数难以计算的问题。

